% Toán tư duy lớp 4 – Chủ đề 1. Ôn tập và bổ sung (bám SGK Toán 4 KNTT, Bài 1–5) – TỰ LUẬN có lời giải từng bước
% Em soạn, đã tự giải lại và kiểm tra mọi phép tính bằng Python (2026-11-18).

% ===================== BÀI 1. ÔN TẬP CÁC SỐ ĐẾN 100 000 =====================

\begin{bt}%[4D1N1-1]
% nguon: Toán tư duy lớp 4 – Bài 1 SGK Toán 4 KNTT
	Số $70\,435$ gồm mấy chục nghìn, mấy nghìn, mấy trăm, mấy chục và mấy đơn vị? Viết số đó thành tổng theo các hàng.
	\loigiai{
		Đọc từ trái sang phải: chữ số $7$ ở hàng chục nghìn, $0$ ở hàng nghìn, $4$ ở hàng trăm, $3$ ở hàng chục, $5$ ở hàng đơn vị.\\
		Vậy số $70\,435$ gồm $7$ chục nghìn, $0$ nghìn, $4$ trăm, $3$ chục và $5$ đơn vị.\\
		$70\,435 = 70\,000 + 400 + 30 + 5$.\\
		Mẹo: hàng nào có chữ số $0$ thì không cần viết vào tổng.
	}
\end{bt}

\begin{bt}%[4D1H1-1]
% nguon: Toán tư duy lớp 4 – Bài 1 SGK Toán 4 KNTT
	Trong số $58\,358$, giá trị của chữ số $5$ ở hàng chục nghìn gấp bao nhiêu lần giá trị của chữ số $5$ ở hàng chục?
	\loigiai{
		Chữ số $5$ ở hàng chục nghìn có giá trị $50\,000$.\\
		Chữ số $5$ ở hàng chục có giá trị $50$.\\
		$50\,000 : 50 = 1\,000$.\\
		Vậy giá trị đó gấp $1\,000$ lần.
	}
\end{bt}

\begin{bt}%[4D1V1-1]
% nguon: Toán tư duy lớp 4 – Bài 1 SGK Toán 4 KNTT
	Tìm số lớn nhất có năm chữ số khác nhau mà tổng các chữ số của nó bằng $20$.
	\loigiai{
		Muốn số lớn nhất thì chữ số ở hàng cao phải lớn nhất có thể.\\
		Hàng chục nghìn chọn $9$, hàng nghìn chọn $8$. Tổng còn lại cho ba chữ số cuối: $20 - 9 - 8 = 3$.\\
		Ba chữ số khác nhau (khác $9$ và $8$) có tổng bằng $3$ chỉ có thể là $2$, $1$, $0$. Xếp giảm dần: $2$, $1$, $0$.\\
		Vậy số cần tìm là $98\,210$ (thử lại: $9+8+2+1+0=20$).
	}
\end{bt}

\begin{bt}%[4D1N1-2]
% nguon: Toán tư duy lớp 4 – Bài 1 SGK Toán 4 KNTT
	Từ ba chữ số $3$, $7$, $0$, hãy viết tất cả các số có ba chữ số khác nhau.
	\loigiai{
		Chữ số hàng trăm phải khác $0$ nên chỉ có thể là $3$ hoặc $7$.\\
		Hàng trăm là $3$: được $307$; $370$.\\
		Hàng trăm là $7$: được $703$; $730$.\\
		Vậy có $4$ số: $307$, $370$, $703$, $730$.
	}
\end{bt}

\begin{bt}%[4D1H1-2]
% nguon: Toán tư duy lớp 4 – Bài 1 SGK Toán 4 KNTT
	Từ năm chữ số $2$, $5$, $8$, $1$, $6$, hãy viết số lớn nhất và số nhỏ nhất có năm chữ số khác nhau. Tính hiệu của hai số đó.
	\loigiai{
		Số lớn nhất: xếp các chữ số từ lớn đến bé: $86\,521$.\\
		Số nhỏ nhất: xếp các chữ số từ bé đến lớn: $12\,568$.\\
		Hiệu: $86\,521 - 12\,568 = 73\,953$.
	}
\end{bt}

\begin{bt}%[4D1V1-2]
% nguon: Toán tư duy lớp 4 – Bài 1 SGK Toán 4 KNTT
	Từ bốn chữ số $1$, $2$, $3$, $4$ có thể lập được bao nhiêu số có bốn chữ số khác nhau? Trong đó có bao nhiêu số chẵn?
	\loigiai{
		Hàng nghìn có $4$ cách chọn, hàng trăm còn $3$ cách, hàng chục còn $2$ cách, hàng đơn vị còn $1$ cách.\\
		Số các số lập được: $4 \times 3 \times 2 \times 1 = 24$ (số).\\
		Số chẵn có chữ số hàng đơn vị là $2$ hoặc $4$.\\
		Nếu hàng đơn vị là $2$: ba chữ số còn lại $1, 3, 4$ xếp được $3 \times 2 \times 1 = 6$ số. Tương tự hàng đơn vị là $4$: $6$ số.\\
		Vậy có $24$ số, trong đó có $6 + 6 = 12$ số chẵn.
	}
\end{bt}

\begin{bt}%[4D1C1-2]
% nguon: Toán tư duy lớp 4 – Bài 1 SGK Toán 4 KNTT
	Có bao nhiêu số có ba chữ số mà tổng các chữ số bằng $5$?
	\loigiai{
		Chữ số hàng trăm khác $0$ và không lớn hơn $5$. Ta xét từng trường hợp hàng trăm:\\
		Hàng trăm $5$: chục + đơn vị $= 0$: số $500$ (1 số).\\
		Hàng trăm $4$: chục + đơn vị $= 1$: $410$, $401$ (2 số).\\
		Hàng trăm $3$: chục + đơn vị $= 2$: $320$, $311$, $302$ (3 số).\\
		Hàng trăm $2$: chục + đơn vị $= 3$: $230$, $221$, $212$, $203$ (4 số).\\
		Hàng trăm $1$: chục + đơn vị $= 4$: $140$, $131$, $122$, $113$, $104$ (5 số).\\
		Tổng cộng: $1 + 2 + 3 + 4 + 5 = 15$ (số).
	}
\end{bt}

\begin{bt}%[4D1N1-3]
% nguon: Toán tư duy lớp 4 – Bài 1 SGK Toán 4 KNTT
	Khi đánh số trang một quyển sách: từ trang $1$ đến trang $9$ cần bao nhiêu chữ số? Từ trang $10$ đến trang $99$ có bao nhiêu trang và cần bao nhiêu chữ số?
	\loigiai{
		Từ trang $1$ đến trang $9$ có $9$ trang, mỗi trang $1$ chữ số: cần $9$ chữ số.\\
		Từ trang $10$ đến trang $99$ có $99 - 10 + 1 = 90$ (trang), mỗi trang $2$ chữ số: cần $90 \times 2 = 180$ (chữ số).
	}
\end{bt}

\begin{bt}%[4D1H1-3]
% nguon: Toán tư duy lớp 4 – Bài 1 SGK Toán 4 KNTT
	Để đánh số trang một quyển sách dày $96$ trang (bắt đầu từ trang $1$) cần dùng bao nhiêu chữ số?
	\loigiai{
		Trang $1$ đến trang $9$: $9$ trang có $1$ chữ số, cần $9$ chữ số.\\
		Trang $10$ đến trang $96$: có $96 - 10 + 1 = 87$ (trang), mỗi trang $2$ chữ số, cần $87 \times 2 = 174$ (chữ số).\\
		Cần tất cả: $9 + 174 = 183$ (chữ số).
	}
\end{bt}

\begin{bt}%[4D1V1-3]
% nguon: Toán tư duy lớp 4 – Bài 1 SGK Toán 4 KNTT
	Một quyển truyện dày $215$ trang. Hỏi phải dùng bao nhiêu chữ số để đánh số trang quyển truyện đó?
	\loigiai{
		Trang có $1$ chữ số (trang $1$ – $9$): $9$ chữ số.\\
		Trang có $2$ chữ số (trang $10$ – $99$): $90$ trang, cần $90 \times 2 = 180$ (chữ số).\\
		Trang có $3$ chữ số (trang $100$ – $215$): $215 - 100 + 1 = 116$ (trang), cần $116 \times 3 = 348$ (chữ số).\\
		Cần tất cả: $9 + 180 + 348 = 537$ (chữ số).
	}
\end{bt}

\begin{bt}%[4D1C1-3]
% nguon: Toán tư duy lớp 4 – Bài 1 SGK Toán 4 KNTT
	Để đánh số trang một quyển sách, người ta đã dùng hết $690$ chữ số. Hỏi quyển sách đó có bao nhiêu trang?
	\loigiai{
		Đánh từ trang $1$ đến trang $99$ cần $9 + 180 = 189$ (chữ số).\\
		Số chữ số dùng cho các trang có ba chữ số: $690 - 189 = 501$ (chữ số).\\
		Số trang có ba chữ số: $501 : 3 = 167$ (trang).\\
		Số trang của quyển sách: $99 + 167 = 266$ (trang).\\
		Thử lại: trang $100$ đến trang $266$ có $167$ trang, cần $501$ chữ số; $189 + 501 = 690$.
	}
\end{bt}

% ===================== BÀI 2. ÔN TẬP CÁC PHÉP TÍNH TRONG PHẠM VI 100 000 =====================

\begin{bt}%[4D1N2-1]
% nguon: Toán tư duy lớp 4 – Bài 2 SGK Toán 4 KNTT
	Tính bằng cách thuận tiện: $2\,345 + 4\,500 + 655 + 5\,500$.
	\loigiai{
		Nhóm các số cộng lại được số tròn:\\
		$2\,345 + 4\,500 + 655 + 5\,500 = (2\,345 + 655) + (4\,500 + 5\,500)$\\
		$= 3\,000 + 10\,000 = 13\,000$.
	}
\end{bt}

\begin{bt}%[4D1H2-1]
% nguon: Toán tư duy lớp 4 – Bài 2 SGK Toán 4 KNTT
	Tính bằng cách thuận tiện:\\
	a) $25 \times 7 \times 4$;\qquad b) $5 \times 36 \times 2$.
	\loigiai{
		Nhớ: $25 \times 4 = 100$ và $5 \times 2 = 10$. Ta đổi chỗ các thừa số để nhân hai số này trước.\\
		a) $25 \times 7 \times 4 = (25 \times 4) \times 7 = 100 \times 7 = 700$.\\
		b) $5 \times 36 \times 2 = (5 \times 2) \times 36 = 10 \times 36 = 360$.
	}
\end{bt}

\begin{bt}%[4D1V2-1]
% nguon: Toán tư duy lớp 4 – Bài 2 SGK Toán 4 KNTT
	Tính nhanh: $1 + 2 + 3 + \dots + 19 + 20$.
	\loigiai{
		Ghép số đầu với số cuối thành từng cặp:\\
		$(1 + 20) + (2 + 19) + (3 + 18) + \dots + (10 + 11)$.\\
		Mỗi cặp có tổng bằng $21$. Có $20$ số nên có $20 : 2 = 10$ cặp.\\
		Tổng: $21 \times 10 = 210$.
	}
\end{bt}

\begin{bt}%[4D1N2-2]
% nguon: Toán tư duy lớp 4 – Bài 2 SGK Toán 4 KNTT
	Tìm các chữ số $a$, $b$ trong phép cộng: $\overline{4a6} + \overline{25b} = 712$.
	\loigiai{
		Hàng đơn vị: $6 + b$ có tận cùng là $2$ nên $6 + b = 12$, suy ra $b = 6$, nhớ $1$.\\
		Hàng chục: $a + 5 + 1$ có tận cùng là $1$ nên $a + 6 = 11$, suy ra $a = 5$, nhớ $1$.\\
		Hàng trăm: $4 + 2 + 1 = 7$, đúng.\\
		Vậy $a = 5$, $b = 6$. Thử lại: $456 + 256 = 712$.
	}
\end{bt}

\begin{bt}%[4D1H2-2]
% nguon: Toán tư duy lớp 4 – Bài 2 SGK Toán 4 KNTT
	Tìm các chữ số $a$, $b$ trong phép trừ: $625 - \overline{1a8} = \overline{b57}$.
	\loigiai{
		Hàng đơn vị: $5$ không trừ được $8$, ta lấy $15 - 8 = 7$, nhớ $1$ sang hàng chục.\\
		Hàng chục: $2$ trừ (nhớ $1$ và $a$) phải ra $5$. Vì $2 < 5$ nên phải mượn: $12 - 1 - a = 5$, suy ra $a = 6$, nhớ $1$ sang hàng trăm.\\
		Hàng trăm: $6 - 1 - 1 = 4$, suy ra $b = 4$.\\
		Vậy $a = 6$, $b = 4$. Thử lại: $625 - 168 = 457$.
	}
\end{bt}

\begin{bt}%[4D1V2-2]
% nguon: Toán tư duy lớp 4 – Bài 2 SGK Toán 4 KNTT
	Tìm các chữ số $a$, $b$, $c$ ($a$ khác $0$) biết: $\overline{abc} + \overline{ab} + a = 271$.
	\loigiai{
		Ta có $\overline{abc} = a \times 100 + b \times 10 + c$ và $\overline{ab} = a \times 10 + b$.\\
		Cộng lại: $a \times 111 + b \times 11 + c = 271$.\\
		Nếu $a = 1$ thì $b \times 11 + c = 160$; nhưng lớn nhất $9 \times 11 + 9 = 108$, không được.\\
		Nếu $a \ge 3$ thì $a \times 111 \ge 333 > 271$, không được. Vậy $a = 2$.\\
		Khi đó $b \times 11 + c = 271 - 222 = 49$. Vì $c \le 9$ nên $b = 4$, $c = 49 - 44 = 5$.\\
		Vậy $a = 2$, $b = 4$, $c = 5$. Thử lại: $245 + 24 + 2 = 271$.
	}
\end{bt}

\begin{bt}%[4D1N2-3]
% nguon: Toán tư duy lớp 4 – Bài 2 SGK Toán 4 KNTT
	Điền dấu $+$ hoặc $-$ vào mỗi ô trống để được phép tính đúng: $9 \ \square\ 8 \ \square\ 7 = 10$.
	\loigiai{
		Thử các cách điền:\\
		$9 + 8 + 7 = 24$; \quad $9 + 8 - 7 = 10$; \quad $9 - 8 + 7 = 8$; \quad $9 - 8 - 7$ không trừ được.\\
		Vậy phép tính đúng là $9 + 8 - 7 = 10$.
	}
\end{bt}

\begin{bt}%[4D1H2-3]
% nguon: Toán tư duy lớp 4 – Bài 2 SGK Toán 4 KNTT
	Điền dấu $+$, $-$ hoặc $\times$ vào mỗi ô trống để được phép tính đúng: $6 \ \square\ 6 \ \square\ 6 = 30$.
	\loigiai{
		Kết quả $30$ khá lớn so với $6 + 6 + 6 = 18$ nên cần có phép nhân.\\
		$6 \times 6 = 36$, mà $36 - 6 = 30$.\\
		Vậy phép tính đúng là $6 \times 6 - 6 = 30$.\\
		Lưu ý: trong biểu thức có nhân và trừ, ta làm phép nhân trước.
	}
\end{bt}

\begin{bt}%[4D1V2-3]
% nguon: Toán tư duy lớp 4 – Bài 2 SGK Toán 4 KNTT
	Mỗi chữ cái thay cho một chữ số, chữ cái khác nhau thay cho chữ số khác nhau. Tìm $A$ và $B$ biết: $\overline{AB} \times 9 = \overline{A0B}$.
	\loigiai{
		$\overline{AB} \times 9 = (A \times 10 + B) \times 9 = A \times 90 + B \times 9$.\\
		$\overline{A0B} = A \times 100 + B$.\\
		Hai số bằng nhau nên $A \times 90 + B \times 9 = A \times 100 + B$.\\
		Bớt đi phần giống nhau ở hai bên: $B \times 8 = A \times 10$.\\
		$B \times 8$ phải là số tròn chục và $B$ là chữ số nên $B = 5$ (khi đó $5 \times 8 = 40$), suy ra $A = 4$.\\
		Vậy $A = 4$, $B = 5$. Thử lại: $45 \times 9 = 405$.
	}
\end{bt}

% ===================== BÀI 3. SỐ CHẴN, SỐ LẺ =====================

\begin{bt}%[4D1N3-1]
% nguon: Toán tư duy lớp 4 – Bài 3 SGK Toán 4 KNTT
	Có bao nhiêu số chẵn có hai chữ số?
	\loigiai{
		Các số chẵn có hai chữ số là $10, 12, 14, \dots, 98$; hai số liền nhau hơn kém nhau $2$ đơn vị.\\
		Số các số: $(98 - 10) : 2 + 1 = 45$ (số).\\
		Công thức đếm: (số cuối $-$ số đầu) $:$ khoảng cách $+ 1$.
	}
\end{bt}

\begin{bt}%[4D1H3-1]
% nguon: Toán tư duy lớp 4 – Bài 3 SGK Toán 4 KNTT
	Từ $1$ đến $100$ có bao nhiêu số lẻ? Tính tổng các số lẻ đó.
	\loigiai{
		Các số lẻ: $1, 3, 5, \dots, 99$. Số các số lẻ: $(99 - 1) : 2 + 1 = 50$ (số).\\
		Ghép cặp đầu – cuối: $(1 + 99) + (3 + 97) + \dots$, mỗi cặp bằng $100$, có $50 : 2 = 25$ cặp.\\
		Tổng: $100 \times 25 = 2\,500$.
	}
\end{bt}

\begin{bt}%[4D1V3-1]
% nguon: Toán tư duy lớp 4 – Bài 3 SGK Toán 4 KNTT
	Bạn An viết liên tiếp các số chẵn từ $2$ đến $200$. Hỏi An đã viết tất cả bao nhiêu chữ số?
	\loigiai{
		Số chẵn có một chữ số: $2, 4, 6, 8$: $4$ số, cần $4$ chữ số.\\
		Số chẵn có hai chữ số: $10$ đến $98$: $(98 - 10) : 2 + 1 = 45$ (số), cần $45 \times 2 = 90$ (chữ số).\\
		Số chẵn có ba chữ số: $100$ đến $200$: $(200 - 100) : 2 + 1 = 51$ (số), cần $51 \times 3 = 153$ (chữ số).\\
		Tất cả: $4 + 90 + 153 = 247$ (chữ số).
	}
\end{bt}

\begin{bt}%[4D1N3-2]
% nguon: Toán tư duy lớp 4 – Bài 3 SGK Toán 4 KNTT
	Không cần tính, cho biết kết quả mỗi phép tính sau là số chẵn hay số lẻ:\\
	a) $135 + 247$;\qquad b) $482 - 159$;\qquad c) $25 \times 46$.
	\loigiai{
		Nhớ: lẻ $+$ lẻ $=$ chẵn; chẵn $-$ lẻ $=$ lẻ; tích có một thừa số chẵn thì là số chẵn.\\
		a) $135$ lẻ, $247$ lẻ nên tổng là số chẵn.\\
		b) $482$ chẵn, $159$ lẻ nên hiệu là số lẻ.\\
		c) $46$ chẵn nên tích là số chẵn.
	}
\end{bt}

\begin{bt}%[4D1H3-2]
% nguon: Toán tư duy lớp 4 – Bài 3 SGK Toán 4 KNTT
	Tổng của $5$ số lẻ bất kì là số chẵn hay số lẻ? Vì sao?
	\loigiai{
		Ghép $4$ số lẻ đầu thành $2$ cặp, mỗi cặp lẻ $+$ lẻ $=$ chẵn, nên tổng $4$ số đó là số chẵn.\\
		Cộng thêm số lẻ thứ năm: chẵn $+$ lẻ $=$ lẻ.\\
		Vậy tổng của $5$ số lẻ là số lẻ.
	}
\end{bt}

\begin{bt}%[4D1V3-2]
% nguon: Toán tư duy lớp 4 – Bài 3 SGK Toán 4 KNTT
	Bạn Bình tính tổng $1 + 2 + 3 + \dots + 15$ được kết quả $121$. Không cần tính lại tổng, hãy giải thích vì sao Bình tính sai.
	\loigiai{
		Từ $1$ đến $15$ có các số lẻ $1, 3, 5, 7, 9, 11, 13, 15$, tức là $8$ số lẻ.\\
		Tổng của $8$ số lẻ (ghép thành $4$ cặp lẻ $+$ lẻ) là số chẵn; cộng thêm các số chẵn vẫn được số chẵn.\\
		Vậy tổng phải là số chẵn, mà $121$ là số lẻ nên Bình tính sai. (Tổng đúng là $120$.)
	}
\end{bt}

\begin{bt}%[4D1H3-3]
% nguon: Toán tư duy lớp 4 – Bài 3 SGK Toán 4 KNTT
	Mẹ mua $12$ quả táo và $15$ quả cam. Mẹ muốn xếp tất cả số quả vào hai đĩa sao cho số quả ở hai đĩa bằng nhau. Mẹ có làm được không? Vì sao?
	\loigiai{
		Tổng số quả: $12 + 15 = 27$ (quả).\\
		Muốn chia đều vào hai đĩa thì tổng số quả phải là số chẵn. $27$ là số lẻ.\\
		Vậy mẹ không làm được.
	}
\end{bt}

\begin{bt}%[4D1V3-3]
% nguon: Toán tư duy lớp 4 – Bài 3 SGK Toán 4 KNTT
	Trên bàn có $7$ cái cốc đều úp miệng xuống. Mỗi lần, bạn được lật đúng $2$ cốc bất kì. Hỏi sau một số lần lật, có thể làm cho cả $7$ cốc đều ngửa lên được không?
	\loigiai{
		Lúc đầu có $0$ cốc ngửa (số chẵn).\\
		Mỗi lần lật $2$ cốc thì số cốc ngửa hoặc tăng $2$, hoặc giảm $2$, hoặc giữ nguyên (lật $1$ cốc úp thành ngửa và $1$ cốc ngửa thành úp).\\
		Vậy số cốc ngửa luôn là số chẵn: $0, 2, 4, 6, \dots$\\
		Muốn cả $7$ cốc ngửa thì số cốc ngửa phải là $7$, là số lẻ. Vậy không thể làm được.
	}
\end{bt}

% ===================== BÀI 4. BIỂU THỨC CHỨA CHỮ =====================

\begin{bt}%[4D1N4-1]
% nguon: Toán tư duy lớp 4 – Bài 4 SGK Toán 4 KNTT
	Tính giá trị của biểu thức $125 + a \times 4$ với $a = 30$.
	\loigiai{
		Thay $a = 30$: $125 + 30 \times 4 = 125 + 120 = 245$.\\
		Nhớ: làm phép nhân trước, phép cộng sau.
	}
\end{bt}

\begin{bt}%[4D1H4-1]
% nguon: Toán tư duy lớp 4 – Bài 4 SGK Toán 4 KNTT
	Tính giá trị của biểu thức $(a + b) \times c$ với $a = 18$, $b = 32$, $c = 7$.
	\loigiai{
		Thay số: $(18 + 32) \times 7 = 50 \times 7 = 350$.\\
		Nhớ: làm phép tính trong ngoặc trước.
	}
\end{bt}

\begin{bt}%[4D1V4-1]
% nguon: Toán tư duy lớp 4 – Bài 4 SGK Toán 4 KNTT
	Cho biểu thức $A = 500 - m \times 8$ ($m$ là số tự nhiên).\\
	a) Tìm $m$ để $A = 300$.\\
	b) Với $m$ bằng bao nhiêu thì $A$ có giá trị nhỏ nhất? Giá trị nhỏ nhất đó là bao nhiêu?
	\loigiai{
		a) $500 - m \times 8 = 300$ nên $m \times 8 = 500 - 300 = 200$, $m = 200 : 8 = 25$.\\
		b) $A$ nhỏ nhất khi $m \times 8$ lớn nhất nhưng không vượt quá $500$.\\
		$500 : 8 = 62$ (dư $4$), nên $m \times 8$ lớn nhất là $62 \times 8 = 496$.\\
		Vậy $m = 62$ và giá trị nhỏ nhất của $A$ là $500 - 496 = 4$.
	}
\end{bt}

\begin{bt}%[4D1N4-2]
% nguon: Toán tư duy lớp 4 – Bài 4 SGK Toán 4 KNTT
	Tìm $x$: $x - 2\,345 = 6\,789$.
	\loigiai{
		$x$ là số bị trừ. Muốn tìm số bị trừ, ta lấy hiệu cộng với số trừ.\\
		$x = 6\,789 + 2\,345 = 9\,134$.
	}
\end{bt}

\begin{bt}%[4D1H4-2]
% nguon: Toán tư duy lớp 4 – Bài 4 SGK Toán 4 KNTT
	Tìm $x$: $x \times 5 + 125 = 1\,000$.
	\loigiai{
		Coi $x \times 5$ là một số hạng: $x \times 5 = 1\,000 - 125 = 875$.\\
		$x = 875 : 5 = 175$.\\
		Thử lại: $175 \times 5 + 125 = 875 + 125 = 1\,000$.
	}
\end{bt}

\begin{bt}%[4D1V4-2]
% nguon: Toán tư duy lớp 4 – Bài 4 SGK Toán 4 KNTT
	Tìm $x$: $(x + 15) \times 4 - 60 = 200$.
	\loigiai{
		Gỡ dần từ ngoài vào trong.\\
		$(x + 15) \times 4 = 200 + 60 = 260$.\\
		$x + 15 = 260 : 4 = 65$.\\
		$x = 65 - 15 = 50$.
	}
\end{bt}

\begin{bt}%[4D1C4-2]
% nguon: Toán tư duy lớp 4 – Bài 4 SGK Toán 4 KNTT
	Tìm $x$: $x + x \times 2 + x \times 3 = 4\,200$.
	\loigiai{
		Viết $x = x \times 1$: $x \times 1 + x \times 2 + x \times 3 = x \times (1 + 2 + 3) = x \times 6$.\\
		$x \times 6 = 4\,200$ nên $x = 4\,200 : 6 = 700$.
	}
\end{bt}

\begin{bt}%[4D1N4-3]
% nguon: Toán tư duy lớp 4 – Bài 4 SGK Toán 4 KNTT
	Một hình vuông có cạnh là $a$. Viết biểu thức tính chu vi $P$ của hình vuông. Tính $P$ khi $a = 15$ cm.
	\loigiai{
		Hình vuông có $4$ cạnh bằng nhau nên $P = a \times 4$.\\
		Với $a = 15$ cm: $P = 15 \times 4 = 60$ (cm).
	}
\end{bt}

\begin{bt}%[4D1H4-3]
% nguon: Toán tư duy lớp 4 – Bài 4 SGK Toán 4 KNTT
	Một hình chữ nhật có chiều dài $a$ và chiều rộng $b$ (cùng đơn vị đo). Viết biểu thức tính chu vi của hình chữ nhật. Tính chu vi và diện tích khi $a = 24$ cm, $b = 16$ cm.
	\loigiai{
		Chu vi: $P = (a + b) \times 2$.\\
		Với $a = 24$ cm, $b = 16$ cm: $P = (24 + 16) \times 2 = 80$ (cm).\\
		Diện tích: $24 \times 16 = 384$ (cm$^2$).
	}
\end{bt}

\begin{bt}%[4D1V4-3]
% nguon: Toán tư duy lớp 4 – Bài 4 SGK Toán 4 KNTT
	Một hình chữ nhật có chiều rộng $a$ cm, chiều dài gấp $3$ lần chiều rộng. Viết biểu thức tính chu vi hình chữ nhật theo $a$. Tính chu vi và diện tích khi $a = 9$.
	\loigiai{
		Chiều dài là $a \times 3$.\\
		Chu vi: $(a + a \times 3) \times 2 = a \times 4 \times 2 = a \times 8$.\\
		Khi $a = 9$: chu vi là $9 \times 8 = 72$ (cm); chiều dài $9 \times 3 = 27$ (cm).\\
		Diện tích: $27 \times 9 = 243$ (cm$^2$).
	}
\end{bt}

% ===================== BÀI 5. GIẢI BÀI TOÁN CÓ BA BƯỚC TÍNH =====================

\begin{bt}%[4D1N5-1]
% nguon: Toán tư duy lớp 4 – Bài 5 SGK Toán 4 KNTT
	Một cửa hàng nhận về $4$ thùng nước mắm, mỗi thùng $24$ chai. Buổi sáng cửa hàng bán $30$ chai, buổi chiều bán $28$ chai. Hỏi cửa hàng còn lại bao nhiêu chai nước mắm?
	\loigiai{
		Bước 1. Số chai nhận về: $24 \times 4 = 96$ (chai).\\
		Bước 2. Số chai đã bán: $30 + 28 = 58$ (chai).\\
		Bước 3. Số chai còn lại: $96 - 58 = 38$ (chai).\\
		Đáp số: $38$ chai.
	}
\end{bt}

\begin{bt}%[4D1H5-1]
% nguon: Toán tư duy lớp 4 – Bài 5 SGK Toán 4 KNTT
	Hà mua $3$ quyển vở, mỗi quyển giá $8\,000$ đồng và $2$ cái bút, mỗi cái giá $6\,500$ đồng. Hà đưa người bán tờ $50\,000$ đồng. Hỏi người bán phải trả lại Hà bao nhiêu tiền?
	\loigiai{
		Bước 1. Tiền mua vở: $8\,000 \times 3 = 24\,000$ (đồng).\\
		Bước 2. Tiền mua bút: $6\,500 \times 2 = 13\,000$ (đồng).\\
		Bước 3. Tiền trả lại: $50\,000 - (24\,000 + 13\,000) = 13\,000$ (đồng).\\
		Đáp số: $13\,000$ đồng.
	}
\end{bt}

\begin{bt}%[4D1V5-1]
% nguon: Toán tư duy lớp 4 – Bài 5 SGK Toán 4 KNTT
	Một trường có khối Ba, khối Bốn và khối Năm. Khối Ba có $156$ học sinh, khối Bốn nhiều hơn khối Ba $24$ học sinh, khối Năm ít hơn khối Bốn $18$ học sinh. Hỏi ba khối có tất cả bao nhiêu học sinh?
	\loigiai{
		Khối Bốn: $156 + 24 = 180$ (học sinh).\\
		Khối Năm: $180 - 18 = 162$ (học sinh).\\
		Cả ba khối: $156 + 180 + 162 = 498$ (học sinh).\\
		Đáp số: $498$ học sinh.
	}
\end{bt}

\begin{bt}%[4D1H5-2]
% nguon: Toán tư duy lớp 4 – Bài 5 SGK Toán 4 KNTT
	Tìm một số, biết rằng lấy số đó cộng với $15$, được bao nhiêu đem nhân với $3$ thì được $120$.
	\loigiai{
		Tính ngược từ cuối, làm phép tính ngược lại (nhân thành chia, cộng thành trừ):\\
		Trước khi nhân với $3$: $120 : 3 = 40$.\\
		Trước khi cộng $15$: $40 - 15 = 25$.\\
		Số cần tìm là $25$. Thử lại: $(25 + 15) \times 3 = 120$.
	}
\end{bt}

\begin{bt}%[4D1V5-2]
% nguon: Toán tư duy lớp 4 – Bài 5 SGK Toán 4 KNTT
	Lan có một số kẹo. Lan cho em $5$ cái, rồi cho bạn một nửa số kẹo còn lại thì Lan còn $8$ cái. Hỏi lúc đầu Lan có bao nhiêu cái kẹo?
	\loigiai{
		Tính ngược từ cuối:\\
		Lan còn $8$ cái là một nửa số kẹo sau khi cho em, nên sau khi cho em Lan có $8 \times 2 = 16$ (cái).\\
		Lúc đầu Lan có: $16 + 5 = 21$ (cái).\\
		Đáp số: $21$ cái kẹo.
	}
\end{bt}

\begin{bt}%[4D1C5-2]
% nguon: Toán tư duy lớp 4 – Bài 5 SGK Toán 4 KNTT
	Một người bán trứng. Lần thứ nhất bán một nửa số trứng và thêm $1$ quả. Lần thứ hai bán một nửa số trứng còn lại và thêm $1$ quả thì vừa hết. Hỏi lúc đầu người đó có bao nhiêu quả trứng?
	\loigiai{
		Tính ngược từ lần bán cuối.\\
		Lần hai: bán một nửa và thêm $1$ quả thì hết, nghĩa là $1$ quả đó chính là nửa còn lại. Trước lần hai có $1 \times 2 = 2$ (quả).\\
		Lần một: sau khi bán một nửa và thêm $1$ quả thì còn $2$ quả, nên nửa số trứng lúc đầu là $2 + 1 = 3$ (quả).\\
		Lúc đầu có $3 \times 2 = 6$ (quả).\\
		Thử lại: bán $3 + 1 = 4$ quả, còn $2$; bán $1 + 1 = 2$ quả, hết.
	}
\end{bt}

\begin{bt}%[4D1H5-3]
% nguon: Toán tư duy lớp 4 – Bài 5 SGK Toán 4 KNTT
	Vừa gà vừa chó có $36$ con, đếm được $100$ cái chân. Hỏi có bao nhiêu con gà, bao nhiêu con chó?
	\loigiai{
		Giả sử cả $36$ con đều là gà thì có $36 \times 2 = 72$ (chân).\\
		Số chân bị thiếu: $100 - 72 = 28$ (chân).\\
		Mỗi con chó có nhiều hơn con gà $2$ chân, nên số chó: $28 : 2 = 14$ (con).\\
		Số gà: $36 - 14 = 22$ (con).\\
		Thử lại: $22 \times 2 + 14 \times 4 = 44 + 56 = 100$ (chân).
	}
\end{bt}

\begin{bt}%[4D1V5-3]
% nguon: Toán tư duy lớp 4 – Bài 5 SGK Toán 4 KNTT
	Minh có $15$ tờ tiền gồm hai loại $2\,000$ đồng và $5\,000$ đồng, tổng cộng là $54\,000$ đồng. Hỏi Minh có mấy tờ mỗi loại?
	\loigiai{
		Giả sử cả $15$ tờ đều là loại $2\,000$ đồng thì có $2\,000 \times 15 = 30\,000$ (đồng).\\
		Số tiền còn thiếu: $54\,000 - 30\,000 = 24\,000$ (đồng).\\
		Mỗi tờ $5\,000$ đồng hơn tờ $2\,000$ đồng là $3\,000$ đồng, nên số tờ $5\,000$ đồng: $24\,000 : 3\,000 = 8$ (tờ).\\
		Số tờ $2\,000$ đồng: $15 - 8 = 7$ (tờ).\\
		Thử lại: $7 \times 2\,000 + 8 \times 5\,000 = 14\,000 + 40\,000 = 54\,000$ (đồng).
	}
\end{bt}

\begin{bt}%[4D1C5-3]
% nguon: Toán tư duy lớp 4 – Bài 5 SGK Toán 4 KNTT
	Một đề thi có $20$ câu. Mỗi câu làm đúng được $5$ điểm, mỗi câu làm sai hoặc bỏ trống bị trừ $2$ điểm. Bạn Hà được $65$ điểm. Hỏi Hà làm đúng bao nhiêu câu?
	\loigiai{
		Giả sử Hà làm đúng cả $20$ câu thì được $5 \times 20 = 100$ (điểm).\\
		Hà bị hụt: $100 - 65 = 35$ (điểm).\\
		Mỗi câu sai thay vì đúng thì mất $5$ điểm và bị trừ thêm $2$ điểm, tức là hụt $5 + 2 = 7$ (điểm).\\
		Số câu sai: $35 : 7 = 5$ (câu). Số câu đúng: $20 - 5 = 15$ (câu).\\
		Thử lại: $15 \times 5 - 5 \times 2 = 75 - 10 = 65$ (điểm).
	}
\end{bt}
